\documentclass[10pt]{article}

\usepackage[utf8]{inputenc}

\usepackage[T1]{fontenc}

\usepackage{amsmath}

\usepackage{amsfonts}

\usepackage{amssymb}

\usepackage[version=4]{mhchem}

\usepackage{stmaryrd}

\usepackage{bbold}



\begin{document}

с л у ч а й); 2) конечной комплексной плоскости $\mathbb{C}$, или сфере Римана с одной выколотой точкой (п а р а б ол и ч е с к и й с л уч а й); 3) единичному кругу $D=$ $=\{z \in \mathbb{C}:|z|<1\}$ на плоскости $\mathbb{C}$, или сфере Римана с разрезом положительной длины (г и п е р б о л и ч eс к и й с л у ч а й). Важный результат состоит в том, что любая Р. п., подобная однолистной, конформно эквивалентна нек-рой канонич. области расширенной комплексной плоскости. В качестве такой канонич. области можно взять всю расширенную плоскость с конечным или бесконечным числом разрезов, параллельных действительной оси, причем нек-рые из әтих разрезов могут вырождаться в точки. Как указано выше, в случае односвязной Р. П. канонич. область либо не имеет ни одного разреза (э л л и п т и ч е с к и й т и II), либо разрез вырождается в точку (п а р а б о л и ч е с к й й т и п), либо разрез имеет положительную длину (г ип е р б о ли ч е с к и й ти п). Все три тина односвязных Р. П. конформно различны, хотя последние два из них топологически эквивалентны. П р о б л е м а т ип а, пока (1983) не получившая полного решения, состоит в разыскании дополнительных условий на односвязную Р. п., при к-рых она принадлежит гиперболическому или параболич. типу (см. [6], [7], [10], [11]).



В общем случае произвольной Р. п. $R$ односвязной Р. п. будет всегда ее универсальная накрывающая поверхность $\hat{R}$, к-рая, следовательно, принадлежит одному из трех указанных типов. Сама Р. п. $R$ считается Р. п. ә лли п тич е с ко го, па а а боличес ког о или г и п е р бол и ч е с к о г о типа в соответствии с тем, какого типа будет ее универсальная накрывающая $\widehat{R}$. Эта классификация Р. п. оправдывается следующими соображениями. Пусть $D$ - одна из трех областей: расширенная комплексная плоскость, конечная комплексная плоскость или единичный круг, а $\Lambda-$ нек-рая группа дробно-линейных отображений области $D$ на себя (автоморфизмов), не имеющая неподвижных точек в $D$. Конформное отображение $w=W(q)$ универсальной накрывающей $\widehat{R}$ на $D$ переводит группу $\widehat{\Lambda}$ преобразований накрытия $\hat{R}$, изоморфоную фундаментальной группе $\pi_{1}(R)$, в нек-рую группу $\Lambda$ автоморфизмов области $D$. При этом $w=W(q)$ можно рассматривать как конформное отображение факторпространства $\hat{R} / \hat{\Lambda}$ на факторпространство $D / \Lambda$, а $\hat{R} / \hat{\Lambda}$ можно отождествить с $R$. Таким образом, $w=W(q)$ можно рассматривать как конформное отображение Р. П. $R$ на факторпространство $\hat{D} / \Lambda$ с нек-рой группой автоморфизмов $\Lambda$, изоморфной фундаментальной группе $\pi_{1}(R)$.



Так как Р. п. $R$ әллиптич. типа обязательно односвязна, то группа $\Lambda$ тривиальна, а значит, такая Р. п. есть обязательно Р. п. функции, обратной рациональной. Односвязная Р. п. параболич. типа обязательно является Р. П. Функции, обратной мероморфной в конечной плоскости. Компактная Р. п. рода $g=0, g=1$ и $g>1$ является соответственно Р. П. әллиштического, параболического и гиперболич. типа.



В связи с конформной әквивалентностью Р. П. возникает также вопрос о строении групшы $\Sigma$ конформных автоморфизмов Р. п. $R$. За нек-рыми простыми исключениями эта группа $\Sigma$ дискретна, а для компактных P. п. рода $g>1$ она конечна (т е о р е м а III в а р ц а). Исключительных случаев, когда группа $\Sigma$ непрерывна, всего семь, а именно (указаны представители соответствующих конформных классов): сфера в әллиптич. случае; сфера с одной или двумя выколотыми точками и тор в параболич. случае; круг, круг с выколотой точкой и кольцо в гинерболич. случае.



Важное значение имеет также п р о б л е м а м о д ул е й Р. П. в различных постановках. Это вопрос о возможном описании разнообразия конформно неэквива- лентных Р. п. того или иного вида. Напр., легко устанавливаются следующие положения. Множество видов конформно неәквивалентных двусвязных Р. п., подобных однолистным (колец), зависит от одного действительного параметра (м о д у л я) $k, 0<k<1$; т. е. два кольца $0<r_{v}<|z|<R_{v}, v=1,2$, конформно эквивалентны тогда и только тогда, когда совпадают отношения их радиусов: $k=\left(r_{1} / R_{1}\right)=\left(r_{2} / R_{2}\right)$. Множество видов конформно неэквивалентных $n$-связных Р. п., подобных однолистным, при $n>2$ зависит от $3 n-6$ действительных параметров. Множество видов конформно неэквивалентных замкнутых Р. п. рода $g \geqslant 1$ при $g=1$ зависит от двух действительных параметров, а при $g>1$ - от $6 g-6$ действительных параметров (см. Римановых поверхностей конборлные классы, а также [3], [12], [13], [15], [16]; относительно поведения функций других классов на Р. П. см. Рилановых поверхностей классибикачия).



Важным аспектом теории Р. п. является свяаь с понятием униформизации. Для многозначной, вообще говоря, аналитич. функции



\[

w=f(z)

\]



ее Р. п. $R_{f}$ дает геометрич. средство униформизации: многозначное соотношение (2) заменяется двумя однозначными представлениями:



\[

w=F(p), z=g(p), p \in R_{f},

\]



выражающими переменные $\boldsymbol{z}$ и $w$ однозначно во всей области существования функции (2) как полной аналитич. Функции. С другой стороны, подход К. Вейершгтрасca (K. Weierstrass) к построению понятия полной аналитич. функции (2) состоит в использовании локального униформизирующего параметра $t$, позволяющего выразить переменные $z$ и $w$ аналитически в виде однозначных аналитич. функций $z=z(t)$ и $w=w(t)$ локально в окрестности нек-рой точки $\left(z_{0}, w_{0}\right), w_{0}=f\left(z_{0}\right)$. Задача униформизации в ее простейшей классич. постановке есть задача синтеза әтих двух идей. Требуется заменить соотношение (2) во всей области его определения двумя аналитич. представлениями $z=z(t), w=w(t)$, где $t-$ униформизирующее комплексное переменное, изменяющееся в нек-рой области на плоскости.



Возможность униформизации в указанной постановке установлена П. Кёбе и независимо А. Пуанкаре почти одновременно в 1907. Если.Р. п. $R_{f}$ функции (2) односвязна или подобна однолистной, то задача униформизации сводится к построению конформного отображения $\varphi: R_{f} \rightarrow D$ Р. п. $R_{f}$ на плоскую область $D$. Представления (3) тогда дают искомую униформизацию:



\[

z=g\left[\varphi^{-1}(t)\right], w=F\left[\varphi^{-1}(t)\right], t \in D .

\]



Конформное отображение $f$ на плоскую область существует для P. п. $R_{f}$, подобных однолистным, и только для них (о б щ а я те о р е м а уни фо р м из а ц и и)



В общем случае произвольного аналитич. соотношения (2) Р. п. $R_{f}$ не является подобной однолистной, но ее универсальная накрывающая $\hat{R}_{f}$ односвязна, и следовательно, существует конформное отображение



\[

\psi: \hat{R}_{f} \longrightarrow D

\]



где $D-$ одна из упоминавшихся областей: $\overline{\mathbb{C}}, \mathbb{C}$ или единичный круг. Функция $w=f(z)$ мероморфна на Р. П. $R_{f}$, а следовательно и на $\hat{R}_{f}$, причем она зависит только от проекции $p=p(q), p \in R_{f}$, точ்ки $q \in \hat{R}_{f}$. Таким образом, получается геометрич. униформизация в виде



\[

\boldsymbol{z}=\boldsymbol{g}[p(q)], w=F[p(q)]

\]





\end{document}